% ====================================================================
% Chapter 1: Graphite anode ICA tail theory
% Scope: discharge/delithiation branch; thermodynamic stage + kinetic tail.
% ====================================================================
\documentclass[11pt,a4paper]{article}
\usepackage[margin=22mm]{geometry}
\usepackage{kotex}
\usepackage{amsmath,amssymb,mathtools,bm}
\usepackage{booktabs,longtable,array}
\usepackage{enumitem}
\usepackage{amsthm}
\usepackage{hyperref}
\usepackage{fancyhdr}
\usepackage{lastpage}
\usepackage{setspace}
\setstretch{1.13}
\setlength{\parskip}{0.45em}\setlength{\parindent}{0pt}\setlength{\headheight}{14pt}
\hypersetup{colorlinks=true, linkcolor=blue, urlcolor=blue, citecolor=blue,
  pdftitle={Graphite anode ICA tail theory - Chapter 1}}
\pagestyle{fancy}\fancyhf{}
\lhead{Graphite anode ICA tail theory --- Ch.1}\rhead{\thepage/\pageref{LastPage}}
\renewcommand{\headrulewidth}{0.4pt}
\newtheorem*{groundingbox}{Grounding}
\newtheorem*{boundbox}{Validity bound}
\newtheorem*{flagbox}{Flagged assumption}
\newtheorem*{keybox}{Keystone}
\newcommand{\dd}{\mathrm{d}}
\newcommand{\eff}{\mathrm{eff}}
\newcommand{\eq}{\mathrm{eq}}
\newcommand{\app}{\mathrm{app}}
\newcommand{\drive}{\mathrm{drive}}
\newcommand{\bg}{\mathrm{bg}}
\newcommand{\cell}{\mathrm{cell}}
\newcommand{\ext}{\mathrm{ext}}
\newcommand{\OCV}{\mathrm{OCV}}
\newcommand{\tot}{\mathrm{tot}}
\newcommand{\tail}{\mathrm{tail}}
\newcommand{\AL}[1]{\textsf{[AL-#1]}}
\title{\textbf{리튬이온전지 흑연 음극 ICA($dQ/dV$) tail 이론 --- Chapter 1}\\[0.4em]
\large Thermodynamic stage, electrode-potential-assisted relaxation barrier, and a first-pass fitting approximation}
\author{}\date{}
\begin{document}
\maketitle

\section*{서: 관찰에서 fitting approximation까지}
\addcontentsline{toc}{section}{서: 관찰에서 fitting approximation까지}
본 장의 출발점은 하나의 실험 관찰이다.
\begin{quote}\itshape
흑연 음극의 ICA(incremental capacity, $dQ/dV$) peak에서 상전이에 의해 생기는 tail은 온도가 낮을수록 길게 늘어지고,
온도가 높을수록 상대적으로 짧게 끝난다. 또한 전위 상태와 C-rate가 바뀌면 peak overlap과 tail 모양도 달라진다.
\end{quote}

이 관찰은 equilibrium thermodynamics만으로 바로 설명되지 않는다. Ideal lattice-gas equilibrium peak는 $RT/F$ scale의
thermal width를 가지므로, 다른 효과가 없으면 낮은 온도에서 더 좁아지는 방향을 준다. 따라서 관측되는 low-temperature
long tail은 equilibrium peak width 자체가 아니라, transition progress가 equilibrium target을 유한한 속도로 따라가는
kinetic lag에서 찾아야 한다. 본 장의 목적은 그 kinetic lag를 electrode potential이 보조하는 effective relaxation barrier와
연결하고, 마지막에는 실제 ICA 데이터에서 바로 사용할 수 있는 simple fitting approximation을 도출하는 것이다.

논리 흐름은 다음과 같다.
\begin{enumerate}[label=\textbf{(\arabic*)},leftmargin=2.2em]
\item Charge conservation이 internal anode potential $V_n$을 정한다. $V_n$은 lookup value가 아니라 internal state와 external charge가 함께 결정하는 implicit variable이다.
\item Equilibrium target $\xi_{\eq,j}(V_n,T)$는 transition이 충분히 느린 scan에서 도달할 목표를 준다.
\item 실제 progress $\xi_j$는 finite rate $k_j$로 target을 따라간다. 이 relaxation rate가 작으면 tail이 길어진다.
\item Electrode-potential drive $\mathcal A_j$는 effective relaxation barrier를 낮추는 bounded approximation으로 들어가며, 그 결과 $k_j$와 tail length $L$이 바뀐다.
\item Barrier distribution은 exponential map을 통해 relaxation-length spectrum이 되고, 관측 tail은 single-mode kernel의 중첩으로 표현된다.
\item Dominant single-mode limit에서는 $Y_{\tail}(q)\approx B\exp[-(q-q_a)/L]$가 되어 Chapter 1만으로 first-pass fitting이 가능하다.
\end{enumerate}

\begin{groundingbox}
모든 equation은 먼저 physical quantity를 말로 정의한 뒤 제시하고, 뒤에서 각 항의 meaning과 unit을 해석한다.
Assumption은 \S\ref{sec:convention}의 ledger에서 established, bounded, flagged로 분리한다. Paper 또는 textbook 근거가
확인되지 않은 항목은 theory identity처럼 쓰지 않는다. 본 장은 discontinuous switch, 임의 cap, 임의 threshold를 physical model로 쓰지 않는다.
\end{groundingbox}

\tableofcontents
\newpage

\section{기호와 단위}\label{sec:notation}
진행률 기호는 $\xi$로 통일한다. 일부 문헌의 $\theta$는 같은 종류의 occupation/progress variable로만 언급하고,
본문에서는 쓰지 않는다. 본 장은 discharge 중 graphite anode가 delithiation되는 branch를 다룬다. External discharge
coordinate는 $q=Q_{\ext}/Q_{\cell}$이고 증가 방향으로 둔다.

\begin{longtable}{p{0.18\textwidth} p{0.12\textwidth} p{0.62\textwidth}}
\toprule 기호 & 단위 & 의미 \\ \midrule \endhead
$q$ & --- & normalized external discharge coordinate, $q=Q_{\ext}/Q_{\cell}$ \\
$Q_{\ext},Q_{\cell}$ & C & external passed charge and reference cell capacity \\
$I$ & A & applied current magnitude는 $|I|$, discharge sign은 별도 $s_I$로 둔다 \\
$V_n$ & V & internal graphite anode potential, charge conservation의 solution \\
$V_{n,\app}$ & V & measured apparent anode potential, $V_n+s_I|I|R_n$ \\
$V_{n,\drive}$ & V & transition rate에 들어가는 drive potential; Chapter 1 기본형은 $V_{n,\drive}=V_n$ \\
$V_{n,\OCV}$ & V & equilibrium limit에서 charge conservation이 주는 anode OCV \\
$R_n$ & $\Omega$ & effective polarization coefficient for apparent-axis shift \\
$j$ & --- & discharge curve에서 분리 가능한 effective transition label \\
$\xi_j,\xi_{\eq,j}$ & --- & branch-local transition progress and its equilibrium target, $0\le\xi\le1$ \\
$U_j,w_j$ & V & effective transition center potential and equilibrium width \\
$Q_{j,\tot}$ & C & transition $j$가 discharge capacity에 주는 positive contribution \\
$Q_{\bg},C_{\bg}$ & C, C/V & unresolved background capacity and $C_{\bg}=\partial Q_{\bg}/\partial V_n$ \\
$\mathcal A_j$ & J/mol & electrode-potential drive, $s_{\phi,j}F(V_{n,\drive}-U_j)$ \\
$s_{\phi,j}$ & --- & discharge branch에서 $V_{n,\drive}>U_j$가 progress를 촉진하도록 정한 sign factor \\
$G$ or $g$ & J/mol & intrinsic relaxation barrier value before potential assistance \\
$\Delta G_{\eff,j}$ & J/mol & effective relaxation barrier, $G-\chi_j\mathcal A_j$ \\
$\chi_j$ & --- & Level-A relaxation-barrier sensitivity; Chapter 1에서는 Butler--Volmer transfer coefficient $\beta_j$와 동일시하지 않는다 \\
$k_j,k_0$ & 1/s & relaxation rate and prefactor \\
$L$ & --- & relaxation length in $q$ coordinate, $L=|I|/(Q_{\cell}k_j)$ \\
$L_\varphi$ & V & voltage-axis tail length, local approximation $L_\varphi\simeq |\dd V/\dd q|_{q_a}L$ \\
$\rho_G(G;T)$ & mol/J & intrinsic barrier probability density \\
$p_L(L)$ & 1 & relaxation-length probability density in $L$ coordinate; $p_L(L)\dd L$ is mode fraction \\
$A_L(L)$ & 1 & relaxation-length amplitude spectrum; $A_L(L)\dd L$ is tail progress amplitude carried by modes in $[L,L+\dd L]$ \\
$a(L)$ & --- & residual/capacity/accessibility amplitude factor converting probability spectrum into amplitude spectrum \\
$\Theta,\Theta_{\eq}$ & --- & capacity-weighted aggregate progress and its equilibrium target \\
$\Theta_{\tail},\Theta_0$ & --- & post-peak tail progress and single-mode residual amplitude at $q_a$ \\
$Q_p$ & C & capacity scale for aggregate progress, typically $Q_p=\sum_jQ_{j,\tot}$ over the local transition group \\
$q_a,x$ & ---, --- & tail-window start and local coordinate $x=q-q_a$ \\
$Y_{\tail}$ & measured ICA unit & baseline-subtracted ICA tail signal used in first-pass fitting \\
\bottomrule
\end{longtable}

단위상 $I$는 $\mathrm{C/s}$이고 $Q_{\cell}$은 C이다. Constant-current segment에서
\begin{equation}
Q_{\ext}(t)=\int_0^t|I|\dd t',\qquad q=\frac{Q_{\ext}}{Q_{\cell}},\qquad
\frac{\dd q}{\dd t}=\frac{|I|}{Q_{\cell}}.
\label{eq:q_time}
\end{equation}

\section{Convention과 assumption ledger}\label{sec:convention}
\textbf{Branch convention.} Graphite lithiation literature often describes stage evolution as lithiation order. This chapter instead labels
transitions by the measured discharge/delithiation branch. Thus $j$ is not a crystallographic stage number. It is a branch-local peak label.
For each $j$, $\xi_j=0$ means that the corresponding discharge peak has not yet progressed, and $\xi_j=1$ means that the branch-local
transition has passed.

\textbf{Barrier convention.} $G$ is the intrinsic relaxation barrier. $\Delta G_{\eff}=G-\chi\mathcal A$ is not a microscopic forward barrier
in the Butler--Volmer sense. It is a Level-A effective barrier controlling the scalar relaxation mobility by which $\xi_j$ approaches
$\xi_{\eq,j}$. Directional forward/backward barrier splitting is reserved for the reaction-kinetic extension in Chapter 3.

\begin{longtable}{p{0.10\textwidth} p{0.16\textwidth} p{0.30\textwidth} p{0.34\textwidth}}
\toprule ID & status & source type & allowed use \\ \midrule \endhead
\AL{1} & established & charge conservation in porous-electrode theory \cite{doyle1993,newman} & $V_n$ is determined by charge balance. \\
\AL{2} & established/bounded & lattice-gas and regular-solution thermodynamics \cite{mckinnon1983,bazant2013} & ideal logistic derivation; nonideal effective width only as bounded coarse-graining. \\
\AL{3} & established & graphite staging observations \cite{dahn1991,ohzuku1993,funabiki1999ea,funabiki1999jes} & staging transitions generate ICA peaks. \\
\AL{4} & bounded & low-rate OCV/GITT calibration & $U_j,w_j,Q_{j,\tot},Q_{\bg}$ are fitted equilibrium-stage quantities. \\
\AL{5} & validity gate & positive differential capacity / implicit-function condition \cite{newman,bazant2013} & $C_{\bg}>0$ is required in the fitting window. \\
\AL{6} & established & ohmic and transport polarization separation \cite{newman,bardfaulkner} & $V_{n,\app}$ and $V_n$ must be separated before interpreting $\chi_j$. \\
\AL{7} & mathematical gate & range and monotonicity condition & charge-balance solution exists only in an admissible window. \\
\AL{10} & established & transition-state theory \cite{eyring1935} & rate scale follows $\exp[-\Delta G^\ddagger/(RT)]$. \\
\AL{11} & bounded & linear free-energy/BEP/BV analogy \cite{evanspolanyi1938,bardfaulkner,bazant2013} & first-order potential assistance of a scalar relaxation barrier near the transition. \\
\AL{12} & established/bounded & Onsager relaxation and two-state kinetic algebra \cite{onsager1931,degrootmazur1962} & local first-order relaxation toward a target. \\
\AL{13} & flagged/bounded & heterogeneous electrode microstructure plus graphite cooperative transition literature & barrier-only independent mode distribution is a mean-field approximation. \\
\AL{14} & established mechanism, bounded application & relaxation-time distributions \cite{svare2000,johnston2006,lindsey1980} & barrier distribution can map to broad relaxation spectrum; graphite ICA use is a model contribution. \\
\AL{15} & established bound & Marcus curvature \cite{marcus1956} & linear barrier assistance is a small-drive approximation. \\
\AL{16} & bounded analogy & Fredholm ratio-substitution methods \cite{lee2011,son2013,jcp2017} & unknown-ratio closure idea only; Fredholm formulas are not copied into the Volterra problem. \\
\bottomrule
\end{longtable}

\section{흑연 staging과 effective transition}\label{sec:staging}
Graphite lithiation/delithiation proceeds through staged intercalation structures, and stage transformations create pronounced ICA peaks
\cite{dahn1991,ohzuku1993,funabiki1999ea,funabiki1999jes}. Because several microscopic stage transformations can overlap in measured
$dQ/dV$, this chapter uses effective transitions $j=1,\dots,N_p$ rather than claiming a one-to-one mapping between every measured peak and
a crystallographic stage boundary.

A single effective transition can still contain many local modes. Local particle size, defect density, grain boundary, stress, and local
accessibility can change the relaxation barrier. In Chapter 1 we first treat this heterogeneity as a distribution of intrinsic barriers $G$ while
holding the equilibrium target of the effective transition fixed. This is the \emph{barrier-only disorder} approximation. It is not a statement
that real graphite phase-boundary motion is microscopically independent; it is a reduced model whose consequence is a relaxation-length
spectrum.

\section{무대 I: charge conservation과 internal potential}\label{sec:charge}
External current fixes how much charge has passed, but it does not directly fix the internal anode potential. The internal potential is the
value that makes stored charge equal to passed charge:
\begin{equation}
\boxed{\;Q_{\cell}q=Q_{\bg}(V_n,T)+\sum_{j=1}^{N_p}Q_{j,\tot}\xi_j.\;}
\label{eq:charge_balance}
\end{equation}
The left-hand side is known from the experiment. The right-hand side depends on the internal potential and transition progress. Given
$(q,T,\{\xi_j\})$, Eq.~\eqref{eq:charge_balance} determines $V_n$ implicitly.

The background term $Q_{\bg}$ is not a visual baseline. It is the capacity contribution not resolved into the selected effective transitions.
Its differential capacity
\begin{equation}
C_{\bg}(V_n,T)=\frac{\partial Q_{\bg}}{\partial V_n}
\label{eq:Cbg}
\end{equation}
must be positive in the local fitting window:
\begin{equation}
C_{\bg}(V_n,T)>0.
\label{eq:monotone}
\end{equation}
This is a validity condition, not a numerical patch. With this monotonicity and with the target charge lying in the range of $Q_{\bg}$, the
implicit solution for $V_n$ is locally unique.

At equilibrium, $\xi_j=\xi_{\eq,j}(V_n,T)$, and the same charge-balance equation gives the equilibrium anode potential
$V_{n,\OCV}(q,T)$. Thus OCV is the equilibrium special solution of charge conservation, not an independent replacement for it.

The measured apparent potential contains polarization:
\begin{equation}
V_{n,\app}=V_n+s_I|I|R_n(q,T,|I|).
\label{eq:Vapp}
\end{equation}
Therefore converting measured data to the internal axis requires both charge-balance consistency and a polarization correction. In a first-pass
fit, $R_n$ should be constrained from low-rate, pulse, current-interrupt, or impedance information before interpreting tail changes as $\chi_j$.

\section{무대 II: equilibrium target}\label{sec:eq}
A minimal equilibrium target comes from a lattice-gas or regular-solution picture. For occupation/progress $\xi$, the chemical potential is
\begin{equation}
\mu(\xi)=\mu^0+RT\ln\frac{\xi}{1-\xi}+\Omega_j(1-2\xi).
\label{eq:mu_latticegas}
\end{equation}
The logarithmic term is the configurational entropy contribution. For $N$ sites and $n=N\xi$ occupied sites,
$W=N!/[n!(N-n)!]$ and Stirling's approximation gives
\begin{equation}
S_{\mathrm{mix}}=-R[\xi\ln\xi+(1-\xi)\ln(1-\xi)]
\end{equation}
per mole of sites. Differentiating $G_{\mathrm{mix}}=-TS_{\mathrm{mix}}$ with respect to $\xi$ gives
$RT\ln[\xi/(1-\xi)]$. The regular-solution enthalpy term $\Omega_j\xi(1-\xi)$ gives $\Omega_j(1-2\xi)$.

In the ideal limit $\Omega_j=0$, equilibrium with electrode potential gives
\begin{equation}
RT\ln\frac{\xi_{\eq,j}}{1-\xi_{\eq,j}}=s_{\phi,j}F(V_n-U_j),
\label{eq:eq_balance}
\end{equation}
where the standard-state chemical potential has been absorbed into $U_j$. Solving step by step, set
$z=s_{\phi,j}(V_n-U_j)/w_j$ and $w_j=RT/F$ in the ideal case. Then
\begin{equation}
\boxed{\;\xi_{\eq,j}(V_n,T)=\frac{1}{1+\exp[-s_{\phi,j}(V_n-U_j)/w_j]}\;}
\label{eq:logistic}
\end{equation}
for the ideal target. Nonzero $\Omega_j$, particle heterogeneity, and finite plateau width can change the shape; then $w_j$ is an effective
width calibrated from low-rate equilibrium data rather than forced to equal $RT/F$.

The derivative of the ideal target is
\begin{equation}
\frac{\partial\xi_{\eq,j}}{\partial V_n}=\frac{s_{\phi,j}}{w_j}\xi_{\eq,j}(1-\xi_{\eq,j}).
\label{eq:dxidV}
\end{equation}
Thus the ideal thermal contribution narrows as $T$ decreases because $RT/F$ decreases. This does not prove that real graphite equilibrium
peaks always narrow exactly as $RT/F$; graphite staging is cooperative and must be calibrated by low-rate OCV/GITT. It does show that the
observed low-temperature long tail should not be attributed to a purely ideal thermal broadening term without additional evidence.

\section{주역 I: electrode-potential-assisted relaxation barrier}\label{sec:rate}
The kinetic variable is not the equilibrium target itself but the rate at which the actual progress approaches it. Define the local electrode
potential drive
\begin{equation}
\mathcal A_j=s_{\phi,j}F(V_{n,\drive}-U_j).
\label{eq:Aj}
\end{equation}
It has unit J/mol. In the discharge branch convention used here, $s_{\phi,j}$ is chosen so that $\mathcal A_j>0$ means the local potential
state assists progress through the selected discharge transition.

Transition-state theory gives a relaxation rate proportional to a Boltzmann factor:
\begin{equation}
k_j\propto \exp\left[-\frac{\Delta G_j^\ddagger}{RT}\right].
\end{equation}
Near the transition center, this chapter uses a bounded linear free-energy approximation for the scalar relaxation barrier:
\begin{equation}
\boxed{\;\Delta G_{\eff,j}=G-\chi_j\mathcal A_j.\;}
\label{eq:Geff}
\end{equation}
Here $G$ is the intrinsic relaxation barrier of a local mode and $\chi_j$ measures how strongly electrode-potential drive changes the
\emph{relaxation mobility}. This $\chi_j$ is not yet the Butler--Volmer transfer coefficient $\beta_j$. Butler--Volmer/Marcus theory normally
splits forward and reverse barriers and therefore changes the forward/reverse ratio. Chapter 1 deliberately stays at Level A: the equilibrium
target is set by thermodynamics, while $\Delta G_{\eff}$ changes how fast the system approaches that target. Directional kinetics are added in
Chapter 3.

With an Eyring-type prefactor,
\begin{equation}
\boxed{\;k_j(T,V_{n,\drive})=k_0(T)\exp\left[-\frac{G-\chi_j\mathcal A_j}{RT}\right],\qquad
k_0(T)=\frac{k_BT}{h}\kappa(T)\exp\left(\frac{\Delta S_a}{R}\right).\;}
\label{eq:kact}
\end{equation}
The exact prefactor may be grouped differently, but the important point for tail shape is the exponential sensitivity of $k_j$ to
$G-\chi_j\mathcal A_j$. If $|\mathcal A_j|$ becomes comparable to the reorganization/curvature scale, the linear barrier approximation is no
longer a global law; it is then a small-drive local approximation.

\begin{keybox}
A useful two-state algebra explains why a first-order relaxation equation can be used without inventing a new target. Let
$r_j^+$ and $r_j^-$ be local forward and reverse rates. Then
\begin{equation}
\dot\xi_j=r_j^+(1-\xi_j)-r_j^-\xi_j=(r_j^++r_j^-)
\left[\frac{r_j^+}{r_j^++r_j^-}-\xi_j\right].
\end{equation}
Thus $k_j=r_j^++r_j^-$ and $\xi_{\eq,j}=r_j^+/(r_j^++r_j^-)$ if the local rates are treated as constants over a small window. Chapter 1
uses $k_j$ as the scalar relaxation rate. It does not claim that $\chi_j\mathcal A_j$ is the microscopic forward barrier lowering of Chapter 3.
\end{keybox}

If another physical process limits the relaxation, a lumped first-order approximation may be written as a serial time-scale sum,
\begin{equation}
\frac{1}{k_j^{\eff}}=\frac{1}{k_j}+\frac{1}{k_{j,\lim}}.
\label{eq:kphys}
\end{equation}
This is a bounded reduced model for transport, finite-site, diffusion, or accessibility bottlenecks. It should not be added as a free fitting knob
without independent evidence.

\section{주역 II: progress dynamics와 single-mode kernel}\label{sec:dyn}
For a local mode in a constant-current segment, the first-order relaxation model is
\begin{equation}
\frac{\dd\xi_j}{\dd t}=k_j[\xi_{\eq,j}(V_n,T)-\xi_j].
\label{eq:dxidt}
\end{equation}
Using Eq.~\eqref{eq:q_time},
\begin{equation}
\frac{\dd\xi_j}{\dd q}=\frac{Q_{\cell}k_j}{|I|}[\xi_{\eq,j}-\xi_j].
\label{eq:dxidq}
\end{equation}

Now examine a post-peak window beginning at $q_a$. Define residual $r=\xi_{\eq,j}-\xi_j$. If the equilibrium target is nearly flat in the
window, $\dd\xi_{\eq,j}/\dd q\simeq0$, and if $\mathcal A_j$ varies slowly enough that $k_j$ is approximately constant, then
\begin{equation}
\frac{\dd r}{\dd q}=-\frac{Q_{\cell}k_j}{|I|}r=-\frac{1}{L}r,
\qquad
\boxed{\;L=\frac{|I|}{Q_{\cell}k_j}.\;}
\label{eq:Ldef}
\end{equation}
The solution is
\begin{equation}
r(q)=r(q_a)\exp\left[-\frac{q-q_a}{L}\right].
\label{eq:residual_decay}
\end{equation}
The observed progress contribution is the derivative, not the residual itself. From Eq.~\eqref{eq:dxidq},
\begin{equation}
\boxed{\;\frac{\dd\xi_j}{\dd q}=\frac{r(q_a)}{L}\exp\left[-\frac{q-q_a}{L}\right].\;}
\label{eq:single_dxidq}
\end{equation}
This $1/L$ factor is the origin of the kernel factor used later. Dimensionally, $L$ is dimensionless because it is measured in the normalized
charge coordinate $q$.

The temperature trend follows immediately. If $T$ decreases and $G-\chi\mathcal A$ remains positive, $k_j$ decreases and $L$ grows.
A larger $L$ means a slower decay in $q$ and therefore a longer tail.

\section{주역 III: barrier distribution에서 relaxation-length spectrum으로}\label{sec:spectrum}
Real electrodes do not have one intrinsic barrier. Let $\rho_G(G;T)$ be a probability density of intrinsic relaxation barriers. In the
barrier-only disorder approximation, modes with different $G$ share the same branch-local equilibrium target but relax at different rates:
\begin{equation}
\xi_j(t)=\int_0^\infty \rho_G(g;T)\xi_j(g,t)\dd g,
\qquad
\frac{\dd\xi_j(g,t)}{\dd t}=k_j(g)[\xi_{\eq,j}-\xi_j(g,t)].
\label{eq:xi_dist}
\end{equation}
This approximation ignores the possibility that local defects also change $U_j,w_j$, and capacity. When those effects dominate, the model
must be extended; Chapter 1 uses this reduced form to isolate the tail effect of barrier heterogeneity.

Insert Eq.~\eqref{eq:kact} into Eq.~\eqref{eq:Ldef}:
\begin{equation}
\boxed{\;L(G,T,V_{n,\drive})=\frac{|I|}{Q_{\cell}k_0(T)}
\exp\left[\frac{G-\chi_j\mathcal A_j}{RT}\right].\;}
\label{eq:L_of_G}
\end{equation}
This is the central temperature amplification. A fixed spread in $G$ becomes a spread in $\ln L$ of size proportional to $1/(RT)$. Therefore
low temperature makes the relaxation-length spectrum broader and gives more weight to long-lived tail modes.

The probability density in length coordinate is obtained by conserving mode fraction:
\begin{equation}
p_L(L)\dd L=\rho_G(G)\dd G.
\end{equation}
The inverse map is $G(L)=\chi_j\mathcal A_j+RT\ln(L/L_0)$ with $L_0=|I|/(Q_{\cell}k_0)$. Since
\begin{equation}
\frac{\dd G}{\dd L}=\frac{RT}{L},
\end{equation}
we get
\begin{equation}
\boxed{\;p_L(L)=\rho_G(G(L);T)\frac{RT}{L},\qquad L\ge L_{\min}=L_0\exp\left[-\frac{\chi_j\mathcal A_j}{RT}\right].\;}
\label{eq:spectrum_prob}
\end{equation}
The lower limit is not an artificial switch. It is the domain condition corresponding to $G\ge0$.

The observed tail needs an amplitude spectrum, not only a probability spectrum. Let $a(L)$ collect the residual amplitude at $q_a$, local
capacity/accessibility, and other independently constrained weights. Then
\begin{equation}
\boxed{\;A_L(L)=a(L)p_L(L)=a(L)\rho_G(G(L);T)\frac{RT}{L},\qquad L\ge L_{\min}.\;}
\label{eq:spectrum}
\end{equation}
$A_L(L)\dd L$ is the amount of post-peak progress amplitude carried by modes whose relaxation length lies in $[L,L+\dd L]$. It is generally
not normalized to one.

\section{관측 tail: kernel integral}\label{sec:kernel}
Each length mode contributes the derivative form in Eq.~\eqref{eq:single_dxidq}. Summing those derivative contributions gives
\begin{equation}
\boxed{\;\frac{\dd\Theta_{\tail}}{\dd q}(q)
\simeq \int_{L_{\min}}^\infty A_L(L)\frac{1}{L}
\exp\left[-\frac{q-q_a}{L}\right]\dd L,\qquad q\ge q_a.\;}
\label{eq:kernel_integral}
\end{equation}
The integration variable is relaxation length. The result is a progress derivative in the external charge coordinate. The factor $RT/L$ inside
$A_L$ came from the $G\to L$ coordinate transformation; the factor $1/L$ in the kernel came from converting residual decay into progress
derivative. These two factors have different origins.

If a single mode dominates, then
\begin{equation}
A_L(L)=\Theta_0\delta(L-L_*),
\end{equation}
and Eq.~\eqref{eq:kernel_integral} becomes
\begin{equation}
\frac{\dd\Theta_{\tail}}{\dd q}\simeq\frac{\Theta_0}{L_*}
\exp\left[-\frac{q-q_a}{L_*}\right].
\label{eq:single_from_kernel}
\end{equation}
Writing only $A_L=\delta(L-L_*)$ would lose the amplitude. The amplitude-bearing delta form is required.

\begin{boundbox}
Barrier-distribution relaxation and stretched relaxation are established in broader relaxation literature \cite{svare2000,johnston2006,lindsey1980}.
Applying that mechanism to graphite ICA tail is a bounded model claim. Therefore this chapter uses $\rho_G$ as a forward-model input or a
constrained family, not as a uniquely invertible quantity extracted from one tail curve.
\end{boundbox}

\section{Self-consistency: where feedback enters}\label{sec:volterra}
Because $\Theta$ contributes to charge storage, it changes $V_n$ through Eq.~\eqref{eq:charge_balance}; because $V_n$ changes the equilibrium
target, the target depends on the unknown progress. This feedback is real, but it is not needed to obtain the first-pass exponential fitting formula.
It is introduced here to show where later corrections enter.

For one mode,
\begin{equation}
\frac{\dd\Theta}{\dd q}=\frac{1}{L}[\Theta_{\eq}(q)-\Theta(q)].
\end{equation}
Multiplying by the integrating factor $e^{q/L}$ and integrating gives
\begin{equation}
\Theta(q)=\Theta(0)e^{-q/L}+\int_0^q\frac{1}{L}e^{-(q-q')/L}\Theta_{\eq}(q')\dd q'.
\label{eq:singlemode_integral}
\end{equation}
Adding spectrum and charge-balance feedback gives a causal Volterra equation,
\begin{equation}
\Theta(q)=\Theta_{\mathrm{init}}(q)+\int_0^q\mathcal K(q,q')
\Theta_{\eq}\left(V_n(q',\Theta(q')),T\right)\dd q',
\label{eq:volterra}
\end{equation}
where
\begin{equation}
\mathcal K(q,q')=\int_{L_{\min}}^\infty A_L(L)\frac{1}{L}e^{-(q-q')/L}\dd L.
\end{equation}
The kernel multiplies the target, not the residual again. Multiplying the kernel by $[\Theta_{\eq}-\Theta]$ at this integral stage would count
the relaxation twice.

\subsection{Plan B: reference evaluation}\label{sec:planB}
Plan B is a reference calculation, not the main theory of Chapter 1. Discretize $G$ into grid points $g_m$, integrate each local ODE, and solve
Eq.~\eqref{eq:charge_balance} for $V_n$ at every time or charge step. This gives a numerical reference for checking approximations when the
spectrum is broad or feedback is strong.

\subsection{Plan A: optional analytic closure}\label{sec:planA}
The user's JCP paper and Refs.~6,7 use an unknown-ratio approximation to solve Fredholm integral equations of the second kind
\cite{lee2011,son2013,jcp2017}. In the 2017 paper, Eq.~(32) is a Fredholm equation and Eq.~(34) replaces an unknown solution ratio by a
reference ratio. The present graphite equation is Volterra, not Fredholm, so the formula is not copied. Only the idea of a controlled ratio
substitution is used.

Linearize the target around a baseline solution $\Theta^{(0)}$:
\begin{equation}
\Theta_{\eq}(V_n(\Theta))\approx a(q)+b(q)\Theta(q),
\qquad
b(q)=\frac{\partial\Theta_{\eq}}{\partial V_n}\frac{\partial V_n}{\partial\Theta}
=-\frac{Q_p}{C_{\bg}}\frac{\partial\Theta_{\eq}}{\partial V_n}.
\label{eq:selfcoupling}
\end{equation}
The sign $b\le0$ holds only in a window where $\partial\Theta_{\eq}/\partial V_n\ge0$ and $C_{\bg}>0$. With the ratio approximation
$\Theta(q')/\Theta(q)\approx\Theta^{(0)}(q')/\Theta^{(0)}(q)$, the closed correction becomes
\begin{equation}
\boxed{\;\Theta_A(q)=\frac{c(q)}{1-\dfrac{1}{\Theta^{(0)}(q)}
\displaystyle\int_0^q\mathcal K(q,q')b(q')\Theta^{(0)}(q')\dd q'}\;,}
\label{eq:closed}
\end{equation}
where $c(q)=\Theta_{\mathrm{init}}(q)+\int_0^q\mathcal K(q,q')a(q')\dd q'$. This closure is bounded: it must be checked against Plan B before
being used as a quantitative fitting expression.

\section{ICA/DVA relation and the simple fitting formula}\label{sec:fiteq}
Differentiate charge conservation with respect to $q$:
\begin{equation}
Q_{\cell}=C_{\bg}\frac{\dd V_n}{\dd q}+\sum_jQ_{j,\tot}\frac{\dd\xi_j}{\dd q}.
\end{equation}
Define $Q_p\Theta=\sum_jQ_{j,\tot}\xi_j$. Then
\begin{equation}
\dd Q=C_{\bg}\dd V_n+Q_p\dd\Theta.
\label{eq:dQsplit}
\end{equation}
Since $\dd Q=Q_{\cell}\dd q$,
\begin{equation}
\boxed{\;\frac{\dd Q}{\dd V_n}=\frac{C_{\bg}(V_n,T)}{1-Q_p(\dd\Theta/\dd Q)}
=\frac{C_{\bg}(V_n,T)}{1-(Q_p/Q_{\cell})(\dd\Theta/\dd q)}.\;}
\label{eq:fiteq}
\end{equation}
Therefore a tail in $\dd\Theta/\dd q$ appears as an ICA tail. But $d\Theta/dq$ being exponential does not mean the measured ICA is exactly
exponential. The exact reduced expression is rational in the tail contribution. Exponential fitting is a small-tail, locally constant-baseline
approximation.

For the apparent axis,
\begin{equation}
\frac{\dd V_{n,\app}}{\dd q}=\frac{\dd V_n}{\dd q}+s_I|I|\frac{\partial R_n}{\partial q},
\qquad
\frac{\dd Q_{\ext}}{\dd V_{n,\app}}=\frac{Q_{\cell}}{\dd V_{n,\app}/\dd q}.
\label{eq:app_axis}
\end{equation}
Thus $L_\varphi\simeq |\dd V_{n,\app}/\dd q|_{q_a}L$. Replacing $V_{n,\app}$ by $V_n$ in this expression is allowed only when
$\partial R_n/\partial q$ is negligible in the chosen tail window.

\subsection{Simple real-data fitting approximation}\label{sec:ch1_simplefit}
Let $x=q-q_a$. In a single-dominant-mode tail,
\begin{equation}
\frac{\dd\Theta_{\tail}}{\dd q}\simeq\frac{\Theta_0}{L}e^{-x/L},
\qquad
\boxed{\;L=\frac{|I|}{Q_{\cell}k_0(T)}
\exp\left[\frac{G-\chi_j\mathcal A_j}{RT}\right].\;}
\label{eq:simplefit}
\end{equation}
Define the baseline-subtracted ICA tail signal by
\begin{equation}
Y_{\tail}(q)=\left(\frac{\dd Q}{\dd V}\right)_{\mathrm{meas}}
-\left(\frac{\dd Q}{\dd V}\right)_{\mathrm{baseline}}.
\label{eq:Ytail_def}
\end{equation}
The baseline is the slowly varying contribution obtained from low-rate equilibrium fit, local polynomial/spline constrained outside the peak,
or the Chapter 1 charge-balance background. It is not a free curve allowed to absorb the tail.

If $(Q_p/Q_{\cell})(\dd\Theta/\dd q)\ll1$ and the voltage-axis conversion is nearly constant over the selected tail window, Eq.~\eqref{eq:fiteq}
reduces to
\begin{equation}
\boxed{\;Y_{\tail}(q)\approx B\exp[-(q-q_a)/L].\;}
\label{eq:Ytail_simple}
\end{equation}
The amplitude $B$ contains $C_{\bg}$, $Q_p/Q_{\cell}$, $\Theta_0/L$, and the chosen voltage-axis conversion. The shape parameter is $L$.

\textbf{First-pass fitting protocol.}
\begin{enumerate}[label=\textbf{F\arabic*.},leftmargin=2.4em]
\item Fix the equilibrium stage from low-rate OCV/GITT: $U_j,w_j,Q_{j,\tot},Q_{\bg}$.
\item Constrain $R_n$ independently before fitting $\chi_j$.
\item Choose a post-peak tail window after the peak maximum where equilibrium derivative is below the measurement noise floor or where the
      semi-log residual is locally linear. The rule must be tied to data noise, not an arbitrary universal number.
\item Fix $q_a$ before fitting. If $q_a$ is free, it becomes strongly correlated with $L$.
\item Fit $\ln Y_{\tail}=\ln B-(q-q_a)/L$ over the selected window. The slope gives $L$.
\item Map $L$ to an effective barrier using
\begin{equation}
G-\chi_j\mathcal A_j=RT\ln\left(\frac{Q_{\cell}k_0(T)L}{|I|}\right).
\label{eq:barrier_from_L}
\end{equation}
When $k_0(T)$ is not independently known, use multi-temperature or multi-potential differences rather than a single absolute value.
\item For Arrhenius analysis, account for the Eyring prefactor and entropy term. A simple slope of $\ln(1/L)$ versus $1/T$ represents
      $-(\Delta H_a-\chi\mathcal A)/R$ only when $\mathcal A$ is held fixed and prefactor variation is either corrected or negligible.
\item At fixed temperature, a local potential-dependence test gives
\begin{equation}
\frac{\partial\ln L}{\partial V_{n,\drive}}=-\frac{\chi_js_{\phi,j}F}{RT}
\label{eq:chi_slope}
\end{equation}
inside the small-drive window.
\end{enumerate}

If the semi-log tail has systematic curvature, the single-mode approximation has failed. That curvature is not merely a fitting inconvenience;
it indicates a relaxation-length spectrum, changing $L(q)$, or self-consistency feedback.

\section{Temperature, potential, and C-rate effects}\label{sec:falsify}
The reduced formula separates three contributions:
\begin{longtable}{p{0.25\textwidth} p{0.30\textwidth} p{0.35\textwidth}}
\toprule layer & terms & role \\ \midrule \endhead
Equilibrium stage & $U_j(T),w_j(T),Q_{\bg}(V_n,T)$ & peak location and equilibrium target \\
Kinetic lag & $k_j(T,V_{n,\drive}),\rho_G(G;T)$ & tail length and spectrum \\
Apparent-axis polarization & $R_n(q,T,|I|)$ & measured voltage-axis shift and distortion \\
\bottomrule
\end{longtable}

Temperature affects the exponential factor in $L$. Lower $T$ amplifies $G/(RT)$ and broadens the distribution in $\ln L$, so low temperature
naturally gives long tails if the kinetic-lag model is valid. Electrode potential affects $L$ through
\begin{equation}
\frac{\partial\ln L}{\partial V_{n,\drive}}=-\frac{\chi_js_{\phi,j}F}{RT}
\end{equation}
for the discharge sign convention. Potential assistance therefore tends to shorten the tail by increasing $k_j$.

C-rate has competing effects. At fixed $k_j$, $L=|I|/(Q_{\cell}k_j)$ increases with current in the $q$ coordinate. But increasing current can
also shift $V_{n,\drive}$ and increase $k_j$ through potential assistance, which decreases $L$. Apparent polarization also shifts and overlaps
peaks on the measured voltage axis. Therefore C-rate trends must be interpreted as competition among current prefactor, potential-assisted
rate increase, and voltage-axis polarization; no unconditional ``higher C-rate means shorter tail'' rule is assumed.

The model is falsifiable by the following checks.
\begin{enumerate}[label=\textbf{N\arabic*.},leftmargin=2.4em]
\item Low-rate equilibrium data must support the chosen equilibrium target shape and $w_j$.
\item The extracted $L$ must show Arrhenius-type behavior only after controlling $\mathcal A_j$ and prefactor variation.
\item The local slope $\partial\ln L/\partial V_{n,\drive}$ must have the sign and scale implied by Eq.~\eqref{eq:chi_slope} in the small-drive window.
\item $\rho_G$ should be constrained by independent physical information or by a low-dimensional forward model; a single tail curve does not uniquely determine it.
\item If $R_n$ and $\chi_j$ are both left free in one dataset, parameter covariance can destroy physical interpretability. They must be separated by multi-rate, pulse, or independent polarization information.
\end{enumerate}

\section{부록: self-consistency evaluation principle}\label{sec:ch1_numeric}
When the single-mode approximation fails, Plan B provides a reference calculation. Choose a barrier grid $\{g_m\}$ covering the relaxation
lengths visible in the experimental tail window. For each $g_m$, integrate
\begin{equation}
\frac{\dd\xi_j(g_m)}{\dd t}=k_j(g_m)[\xi_{\eq,j}(V_n,T)-\xi_j(g_m)],
\end{equation}
and solve charge conservation for $V_n$ at each step. The algebraic constraint is well posed only when the local charge-balance derivative
condition is satisfied. Plan A is accepted only if its deviation from this reference is below the measurement noise and discretization floor in the
same temperature, rate, and tail window.

\section*{Chapter 2--5로의 전달}
Chapter 1 fixes the convention for $V_n$, $V_{n,\app}$, $V_{n,\drive}$, $\xi_j$, $\xi_{\eq,j}$, $k_j$, $L$, and $A_L$.
Chapter 2 should use the equilibrium stage to discuss reversible heat and entropy coefficients. Chapter 3 should expand Level-A scalar
relaxation into directional forward/backward kinetics. Chapter 4 should combine heat terms without double counting, and Chapter 5 should
add charge/discharge branch hysteresis without changing the Chapter 1 convention.

\begin{thebibliography}{99}
\bibitem{doyle1993} M. Doyle, T. F. Fuller, J. Newman, ``Modeling of Galvanostatic Charge and Discharge of the Lithium/Polymer/Insertion Cell,''
  \emph{J. Electrochem. Soc.} \textbf{140}, 1526 (1993). DOI: 10.1149/1.2221597.
\bibitem{newman} J. Newman, K. E. Thomas-Alyea, \emph{Electrochemical Systems}, 3rd ed., Wiley (2004). ISBN 978-0-471-47756-3.
\bibitem{mckinnon1983} W. R. McKinnon, R. R. Haering, ``Physical Mechanisms of Intercalation,'' in \emph{Modern Aspects of Electrochemistry} No.~15, Plenum, New York (1983), p.~235.
\bibitem{bazant2013} M. Z. Bazant, ``Theory of Chemical Kinetics and Charge Transfer Based on Nonequilibrium Thermodynamics,'' \emph{Acc. Chem. Res.} \textbf{46}, 1144 (2013). DOI: 10.1021/ar300145c.
\bibitem{bardfaulkner} A. J. Bard, L. R. Faulkner, \emph{Electrochemical Methods}, 2nd ed., Wiley (2001). ISBN 978-0-471-04372-0.
\bibitem{eyring1935} H. Eyring, ``The Activated Complex in Chemical Reactions,'' \emph{J. Chem. Phys.} \textbf{3}, 107 (1935). DOI: 10.1063/1.1749604.
\bibitem{evanspolanyi1938} M. G. Evans, M. Polanyi, ``Inertia and Driving Force of Chemical Reactions,'' \emph{Trans. Faraday Soc.} \textbf{34}, 11 (1938). DOI: 10.1039/TF9383400011.
\bibitem{marcus1956} R. A. Marcus, ``On the Theory of Oxidation-Reduction Reactions Involving Electron Transfer. I,'' \emph{J. Chem. Phys.} \textbf{24}, 966 (1956). DOI: 10.1063/1.1742723.
\bibitem{onsager1931} L. Onsager, ``Reciprocal Relations in Irreversible Processes. I,'' \emph{Phys. Rev.} \textbf{37}, 405 (1931). DOI: 10.1103/PhysRev.37.405.
\bibitem{degrootmazur1962} S. R. de Groot, P. Mazur, \emph{Non-Equilibrium Thermodynamics}, North-Holland (1962); Dover (1984).
\bibitem{dahn1991} J. R. Dahn, ``Phase Diagram of Li$_x$C$_6$,'' \emph{Phys. Rev. B} \textbf{44}, 9170 (1991). DOI: 10.1103/PhysRevB.44.9170.
\bibitem{ohzuku1993} T. Ohzuku, Y. Iwakoshi, K. Sawai, ``Formation of Lithium-Graphite Intercalation Compounds in Nonaqueous Electrolytes,'' \emph{J. Electrochem. Soc.} \textbf{140}, 2490 (1993). DOI: 10.1149/1.2220849.
\bibitem{funabiki1999ea} A. Funabiki, M. Inaba, T. Abe, Z. Ogumi, ``Nucleation and Phase-Boundary Movement upon Stage Transformation in Lithium-Graphite Intercalation Compounds,'' \emph{Electrochim. Acta} \textbf{45}, 865 (1999). DOI: 10.1016/S0013-4686(99)00290-X.
\bibitem{funabiki1999jes} A. Funabiki, M. Inaba, T. Abe, Z. Ogumi, ``Stage Transformation of Lithium-Graphite Intercalation Compounds Caused by Electrochemical Lithium Intercalation,'' \emph{J. Electrochem. Soc.} \textbf{146}, 2443 (1999). DOI: 10.1149/1.1391953.
\bibitem{baessler1993} H. B\"assler, ``Charge Transport in Disordered Organic Photoconductors,'' \emph{Phys. Status Solidi B} \textbf{175}, 15 (1993). DOI: 10.1002/pssb.2221750102.
\bibitem{svare2000} I. Svare, S. W. Martin, F. Borsa, ``Stretched Exponentials with $T$-Dependent Exponents from Fixed Distributions of Energy Barriers for Relaxation Times in Fast-Ion Conductors,'' \emph{Phys. Rev. B} \textbf{61}, 228 (2000). DOI: 10.1103/PhysRevB.61.228.
\bibitem{johnston2006} D. C. Johnston, ``Stretched Exponential Relaxation Arising from a Continuous Sum of Exponential Decays,'' \emph{Phys. Rev. B} \textbf{74}, 184430 (2006). DOI: 10.1103/PhysRevB.74.184430.
\bibitem{lindsey1980} C. P. Lindsey, G. D. Patterson, ``Detailed Comparison of the Williams-Watts and Cole-Davidson Functions,'' \emph{J. Chem. Phys.} \textbf{73}, 3348 (1980). DOI: 10.1063/1.440530.
\bibitem{jcp2017} K. Lee, S. Lee, C. H. Choi, S. Lee, ``Effects of External Electric Field and Anisotropic Long-Range Reactivity on Charge Separation Probability,'' \emph{J. Chem. Phys.} \textbf{147}, 144111 (2017). DOI: 10.1063/1.5000882.
\bibitem{lee2011} S. Lee, C. Y. Son, J. Sung, S.-H. Chong, ``Communication: Propagator for Diffusive Dynamics of an Interacting Molecular Pair,'' \emph{J. Chem. Phys.} \textbf{134}, 121102 (2011). DOI: 10.1063/1.3565476.
\bibitem{son2013} C. Y. Son, J. Kim, J.-H. Kim, J. S. Kim, S. Lee, ``An Accurate Expression for the Rates of Diffusion-Influenced Bimolecular Reactions with Long-Range Reactivity,'' \emph{J. Chem. Phys.} \textbf{138}, 164123 (2013). DOI: 10.1063/1.4802584.
\end{thebibliography}

\end{document}
